% Generated draft: requires mathematical and source review.
\documentclass[11pt]{article}
\usepackage[margin=25mm]{geometry}
\usepackage{fontspec}
\setmainfont{Latin Modern Roman}
\usepackage{amsmath,amssymb,mathtools,unicode-math}
\setmathfont{Latin Modern Math}
\usepackage{enumitem,needspace,xurl,hyperref}
\setlength{\parindent}{0pt}
\setlength{\parskip}{5pt}
\setlength{\emergencystretch}{3em}
\newcommand{\N}{\mathbb{N}}
\newcommand{\Z}{\mathbb{Z}}
\newcommand{\Q}{\mathbb{Q}}
\newcommand{\R}{\mathbb{R}}
\newcommand{\C}{\mathbb{C}}
\newcommand{\F}{\mathbb{F}}
\newcommand{\A}{\mathbb{A}}
\newcommand{\PP}{\mathbb{P}}
\newcommand{\E}{\mathbb{E}}
\newcommand{\eps}{\varepsilon}
\newcommand{\dd}{\,\mathrm{d}}
\begin{document}
\section*{UnsolvedMath ID: 30005207}
\textbf{Source problem code:} catalogue entry
\subsection{Definitions and mathematical statement}
Let \(\Omega\subset\R^n\) be a Lipschitz domain: a connected open set whose boundary is locally representable, after a rigid change of coordinates, as the graph of a Lipschitz function, with \(\Omega\) locally on one side of that graph. No smallness assumption is imposed on the Lipschitz constants. Let \(B=B(x_0,r)\) be a Euclidean ball with \(x_0\in\partial\Omega\) and \(r>0\), and set \(\Sigma=B\cap\partial\Omega\). Let \(u\in C(\overline\Omega)\) be harmonic in \(\Omega\), meaning that \(\Delta u=0\) in the interior. Write \(\mathcal H^{n-1}\) for \((n-1)\)-dimensional Hausdorff measure.

The supplied question asks whether, for every such \(\Omega,B,u\),
\[
 u|_{\Sigma}=0
 \quad\text{and}\quad
 \mathcal H^{n-1}\bigl(\{x\in\Sigma:\nabla u(x)=0\}\bigr)>0
 \quad\Longrightarrow\quad
 u\equiv0\ \text{in }\Omega.
\]
Here \(u|_{\Sigma}=0\) is understood pointwise using continuity on \(\overline\Omega\).

\textbf{Unresolved formulation qualification.} The supplied hypotheses do not define \(\nabla u\) on \(\Sigma\): interior harmonicity and continuity on the closure do not themselves provide a classical boundary gradient at every boundary point. The source record does not specify whether the displayed boundary gradient denotes a nontangential limit of the interior gradient, another trace, or a gradient supplied by additional boundary regularity. The implication above therefore retains the supplied notation as an unresolved part of the question; a precise boundary interpretation and any hypotheses needed for it require mathematical review. The record also does not specify the intended range of \(n\) or whether \(u\) is real-valued or complex-valued.
\subsection{Short English statement}
Does a harmonic function on an arbitrary Lipschitz domain vanish identically if it vanishes on a boundary patch and its boundary gradient vanishes on a subset of positive surface measure? The question removes every smallness assumption on the Lipschitz constant. The meaning of the boundary gradient remains unspecified in the supplied record.
\subsection{Sources}
\begin{itemize}
\item Local UnsolvedMath registry: \texttt{data/unsolvedmath/problems.json}; permanent UnsolvedMath ID \(30005207\); exact source problem code \texttt{catalogue entry}. Title: Boundary Unique Continuation on General Lipschitz Domains.
\item Supplied source citation: \emph{Calculus of Variations (2023)}, Oberwolfach Reports, DOI \texttt{10.4171/owr/2022/37}, pp.\ 2129--2193. \url{https://doi.org/10.4171/owr/2022/37}. The collection describes these reports as records of mathematical workshops and their open-problem discussions.
\item Supplied original excerpt: the question asks whether ``the same statement'' holds for Lipschitz domains after removing the smallness assumption on the Lipschitz constant. The excerpt associates the discussion with Zihui Zhao, \emph{Quantitative boundary unique continuation and the estimate of the singular set}, joint work with Carlos Kenig. No separate proposer is identified; the registry gives proposed year 2022.
\item Supplied literature reference: Tolsa, boundary unique continuation. \url{https://doi.org/10.1002/cpa.22025}.
\item Supplied literature reference: Kenig--Zhao, \(C^1\)-Dini domains. \url{https://arxiv.org/abs/2102.07281}.
\item Supplied literature reference: original Oberwolfach report. \url{https://doi.org/10.4171/owr/2022/37}.
\item Literature assessment dated exactly \texttt{2026-08-22}: the supplied curation labels the problem open and states that boundary unique continuation is proved for \(C^1\)-Dini domains and for Lipschitz domains with sufficiently small constant, while the unrestricted Lipschitz-domain version remains open. This is supplied dated curation derived from the OWR record, not an independently verified literature assessment; its status and cited claims require independent verification.
\item The registry labels the statement \texttt{corrected\_verified} and reports that the official report was checked, the question isolated from surrounding exposition, and notation repaired. These are supplied verification claims, not verification performed for this draft.
\end{itemize}
\end{document}
